\clearpage
\appendix
\section*{附录 A\quad 写作说明与待替换清单（投稿前删除）}
\addcontentsline{toc}{section}{附录 A 写作说明与待替换清单}
\markboth{附录 A\quad 写作说明}{附录 A\quad 写作说明}
\small

\subsection*{A.1 数值的三种状态}
全部数值集中在 \texttt{data/values.yaml}，每项为 \texttt{[值, 状态, 来源]}。\texttt{real}（\NumReal{} 项）来自已有运行，正常显示；\texttt{ideal}（\NumIdeal{} 项）为理想示意值，正文中显示为\textcolor{red}{红色}，相关图片带“示意数据”水印；\texttt{todo}（\NumTodo{} 项）为待确认信息，显示为\textcolor{orange!85!black}{[橙色方括号]}。实验完成后，把实测值填入对应项并把状态改为 \texttt{real}，运行 \texttt{./build.sh} 即可同步更新正文、表格和图片。表~\ref{tab:coverage} 和第~\ref{sec:res_general} 节的数字是真实结果，直接引自英文稿 v0.4（\texttt{paper/results/memory\_coverage.json} 与 RoboMME 评测），未经 \texttt{values.yaml}。

\textbf{示意值不是预测。}示意值是按“实验取得预期结果”设定的、内部自洽的数字（例如失败分类之和等于总次数，分层成功率与总成功率一致），只用于搭建叙事。实测结果的方向或大小与示意值不同时，不仅要改数字，还要改依赖结论方向的句子，主要在：摘要与引言贡献列表的结果句；第~\ref{sec:res_closed}、\ref{sec:res_long} 节各段的小标题和首句；第~\ref{sec:res_ablation} 节“事件来源”“事件扰动训练”两段；第~\ref{sec:res_real} 节；结论。图~\ref{fig:capacity} 与第~\ref{sec:coverage}、\ref{sec:hem} 节的容量数字（约 \V{cap_u}、约 \V{cap_k8}、\V{cap_hem}、\V{cap_flat_max}）是按设计参数真实计算的，由 \texttt{tools/make\_figures\_zh.py} 中的 \texttt{fig\_capacity} 生成；改动 $m_w$、$c$、$s_{\max}$ 后需重新运行。

\subsection*{A.2 尚待实现与完成的实验（对应章节）}
\begin{enumerate}
\item \textbf{R0 捷径审计}（表~\ref{tab:audit}）：交接文档修订 7 第 15.4 节，正在执行；FruitRow 一行需在 R0 之后追加。
\item \textbf{\hem{} 的实现}（第~\ref{sec:hem} 节）：分层选帧函数（含场景锚、自适应分配）、拼图帧构造（脱离后第 3 帧前置图像、以 TCP 投影为中心裁剪 $112\times112$、$2\times2$ 拼接、空格屏蔽）、类型-时间嵌入（零初始化，接到记忆分支 token 上）。训练与部署须调用同一函数；需新增“记忆帧与声明规则一致”的逐决策核验。
\item \textbf{\pedet{} 的实现与训练}（第~\ref{sec:pedet} 节，表~\ref{tab:detect}，图~\ref{fig:events}(c)）：从仿真读取指尖接触与果柄约束生成逐帧标签；采集基础策略在训练布局上的闭环回放；实现域随机化。网络结构与超参数为设计值，按实际实现修改。
\item \textbf{事件扰动训练 EPT}（第~\ref{sec:ept} 节）：在训练样本的记忆构造前对事件表做删除、误插和抖动；扰动率以真机检测误差为准。
\item \textbf{FruitRow 环境}（第~\ref{sec:tasks} 节）：3 株植物、换株信号、整行配额、风落果；专家示范与审计。
\item \textbf{v2 示范、K0/K4/K8/HEM 训练与正式评测}（表~\ref{tab:closed}、图~\ref{fig:simmain}、正对照）：修订 7 第 15.5 节需相应扩展为四个条件和三个任务，主假设族改为 H1--H3；协议须事先写明部署时事件由 \pedet{} 给出、K4/K8 与 HEM 共用 EPT。
\item \textbf{消融}（表~\ref{tab:ablation}）：(a) 为推理期变体，(b) 中各项及“去掉 EPT”需重新训练。
\item \textbf{仿真光照扰动}：渲染亮度增益 0.5--1.6，用于图~\ref{fig:events}(c) 与消融。
\item \textbf{真机实验}（第~\ref{sec:real_design} 节全部）：番茄实验架与弹簧果梗夹（拉脱力标定）、舵机拨杆与叶帘、窄带照明可分性验证、第三视角相机与双人标注流程、遥操作示范、事后确认自训练（持物段阈值、篮区标定、指间区域像素差阈值）、检测评测与闭环评测、温室验证。表~\ref{tab:hardware} 的硬件型号均待确认；图~\ref{fig:realsetup} 待替换为实物照片。
\item \textbf{运行开销}：测量仿真与真机上的 VLA 推理、\pedet{} 和 \hem{} 构造耗时。
\item \textbf{可选}：用 v2 示范重算第~\ref{sec:res_memory} 节的离线覆盖率分析，并加入 \hem{} 一行，使离线与闭环结果来自同一版本的数据。
\end{enumerate}

\subsection*{A.3 需要你决定的事项}
\begin{enumerate}
\item \textbf{方法范围。}\hem{}、拼图帧、类型-时间嵌入、EPT、HSL 和 FruitRow 都是本稿新增的设计，尚未写入交接文档，也未经封存协议。按约定，修改记忆方法与系统定义需要你确认后才能交给执行端实现。
\item \textbf{事件定义。}本稿以“与成熟度无关的物理采摘事件（接触 + 脱离）”为主定义，几何距离事件作为基线。当前代码（\texttt{fruit\_distance.py}）计算的是到“成熟且仍挂树”果实的距离，与本稿不一致。
\item \textbf{基础模型的引用。}正文按要求不突出基础模型的名称，只在第~\ref{sec:base} 节以“当前主流结构”描述并引用 $\pi_0$/$\pi_{0.5}$；记忆分支的结构来源建议仍保留一处引用（第~\ref{sec:related} 节已有），以符合学术规范。
\item \textbf{是否保留 v1 的零结果}（第~\ref{sec:v1} 节与表~\ref{tab:v1}）。这些是真实结果，写成“基准有效性”的动机能增强可信度。
\item \textbf{目标期刊与格式。}本稿按中文期刊常见结构撰写；若投英文期刊，可在定稿后翻译并改用 elsarticle 模板。
\item \textbf{真机规模。}每条件 30 次只能可靠检测约 30~\pp{} 量级的差异，真机 HEM$-$K8 的差距预计不显著；若希望真机上也体现分层优势，需要增加一个真机长时任务（如两株番茄架连续作业）或增加布局数。
\end{enumerate}

\subsection*{A.4 需要核对的参考文献与数字}
\texttt{shi2025memoryvla}、\texttt{fang2025sam2act}、\texttt{kim2024imagegrid} 的作者、题目和出处需要核对；引言中 Bac 等综述的“平均成功率约 66\%、作业周期约 33~s”需对照原文核实；番茄果梗分离力需补充文献。其余文献沿用英文稿 v0.4 的已核对条目，新增的经典文献（RT-1、Q-attention、TCN、域随机化、捷径学习、ResNet、focal loss、伪标签、Noisy Student、HER、bootstrap、Wilson 区间）建议投稿前统一再核一次页码与卷期。

\subsection*{A.5 示意值与待确认信息清单}
{\footnotesize
\begin{longtable}{p{3.2cm}p{0.9cm}p{4.2cm}p{6.2cm}}
\toprule
键名 & 状态 & 当前值 & 说明\\
\midrule
\endhead
\texttt{author1} & todo & 作者一 & \\
\texttt{author2} & todo & 作者二 & \\
\texttt{author3} & todo & 通信作者 & \\
\texttt{affil} & todo & 单位，城市 邮编 & \\
\texttt{email} & todo & email@example.com & \\
\texttt{fund} & todo & 基金项目名称（编号） & \\
\texttt{aud\_cnt\_o1} & ideal & 14 & 计数 v2：仅当前帧脚本策略成功率 \%\\
\texttt{aud\_cnt\_o2} & ideal & 18 & 计数 v2：第 0 帧+当前帧\\
\texttt{aud\_cnt\_o3} & ideal & 96 & 计数 v2：完整历史\\
\texttt{aud\_occ\_o1} & ideal & 26 & 遮挡 v2：仅当前帧\\
\texttt{aud\_occ\_o2} & ideal & 94 & 遮挡 v2：第 0 帧+当前帧\\
\texttt{aud\_occ\_o3} & ideal & 96 & 遮挡 v2：完整历史\\
\texttt{exp\_cnt} & ideal & 93.5 & 计数 v2 专家逐尝试成功率 \%\\
\texttt{exp\_occ} & ideal & 95.0 & 遮挡 v2 专家逐尝试成功率 \%\\
\texttt{clip\_cnt} & ideal & 0.21 & 示范通过安全包络的裁剪率 \%\\
\texttt{clip\_occ} & ideal & 0.18 & \\
\texttt{det\_n\_ep} & ideal & 300 & 检测评测回合数（专家示范 100 + 策略闭环回放 200）\\
\texttt{det\_n\_ev} & ideal & 812 & 真值采摘事件数\\
\texttt{det\_n\_fail} & ideal & 164 & 失败抓取尝试数\\
\texttt{det\_geo\_p} & ideal & 0.832 & \\
\texttt{det\_geo\_r} & ideal & 1.000 & \\
\texttt{det\_geo\_f1} & ideal & 0.908 & \\
\texttt{det\_geo\_cnt} & ideal & 61.3 & 计数正确回合 \%\\
\texttt{det\_geo\_fp} & ideal & 100 & 失败尝试被计为采摘的比例 \%（几何定义下必然）\\
\texttt{det\_geo\_light} & ideal & 1.000 & 光照扰动下 F1（不依赖图像）\\
\texttt{det\_fdm\_p} & ideal & 0.903 & \\
\texttt{det\_fdm\_r} & ideal & 0.981 & \\
\texttt{det\_fdm\_f1} & ideal & 0.940 & \\
\texttt{det\_fdm\_cnt} & ideal & 74.7 & \\
\texttt{det\_fdm\_fp} & ideal & 38 & \\
\texttt{det\_fdm\_light} & ideal & 0.214 & \\
\texttt{det\_grip\_p} & ideal & 0.887 & \\
\texttt{det\_grip\_r} & ideal & 0.952 & \\
\texttt{det\_grip\_f1} & ideal & 0.918 & \\
\texttt{det\_grip\_cnt} & ideal & 66.0 & \\
\texttt{det\_grip\_fp} & ideal & 41 & \\
\texttt{det\_grip\_light} & ideal & 0.918 & \\
\texttt{det\_pv\_p} & ideal & 0.942 & \\
\texttt{det\_pv\_r} & ideal & 0.957 & \\
\texttt{det\_pv\_f1} & ideal & 0.949 & \\
\texttt{det\_pv\_cnt} & ideal & 81.3 & \\
\texttt{det\_pv\_fp} & ideal & 12 & \\
\texttt{det\_pv\_light} & ideal & 0.902 & \\
\texttt{det\_pp\_p} & ideal & 0.901 & \\
\texttt{det\_pp\_r} & ideal & 0.963 & \\
\texttt{det\_pp\_f1} & ideal & 0.931 & \\
\texttt{det\_pp\_cnt} & ideal & 72.0 & \\
\texttt{det\_pp\_fp} & ideal & 33 & \\
\texttt{det\_pp\_light} & ideal & 0.931 & \\
\texttt{det\_pe\_p} & ideal & 0.978 & \\
\texttt{det\_pe\_r} & ideal & 0.983 & \\
\texttt{det\_pe\_f1} & ideal & 0.980 & \\
\texttt{det\_pe\_cnt} & ideal & 92.7 & \\
\texttt{det\_pe\_fp} & ideal & 4 & \\
\texttt{det\_pe\_light} & ideal & 0.962 & \\
\texttt{det\_pe\_tc} & ideal & 1 & 接触关键帧误差中位数（帧）\\
\texttt{det\_pe\_tr} & ideal & 0 & 脱离关键帧误差中位数（帧）\\
\texttt{det\_pe\_tr90} & ideal & 2 & 脱离关键帧误差 P90（帧）\\
\texttt{det\_pe\_kf} & ideal & 0.987 & K8 记忆关键帧召回（±3 帧）\\
\texttt{pedet\_params} & ideal & 1.9 & 时序与融合模块参数量（M），以实际实现为准\\
\texttt{pedet\_ms} & ideal & 1.6 & 检测头单帧延迟 ms（L20）\\
\texttt{fig\_light\_pedet} & ideal & 0.951, 0.972, 0.98, 0.975, 0.958 & \\
\texttt{fig\_light\_fdm} & ideal & 0.12, 0.48, 0.94, 0.61, 0.15 & \\
\texttt{fig\_light\_grip} & ideal & 0.918, 0.918, 0.918, 0.918, 0.918 & \\
\texttt{hem\_params} & ideal & 0.14 & 类型与时间嵌入新增参数（M），按实现修改\\
\texttt{hem\_ms} & ideal & 0.21 & HEM 选帧与拼图构建耗时 ms（CPU）\\
\texttt{fm\_steps} & todo & 10 & 流匹配积分步数，按实际配置确认\\
\texttt{c\_k0} & ideal & 58, 45, 22, 9, 16 & K0 均匀记忆\\
\texttt{c\_k4} & ideal & 78, 30, 16, 10, 16 & 平坦 PEKM k=4\\
\texttt{c\_k8} & ideal & 88, 20, 12, 10, 20 & 平坦 PEKM k=8\\
\texttt{c\_hem} & ideal & 96, 14, 10, 10, 20 & HEM（完整 HarvestMem）\\
\texttt{c\_k0\_s} & ideal & 38.7 & 成功率 \%\\
\texttt{c\_k4\_s} & ideal & 52.0 & \\
\texttt{c\_k8\_s} & ideal & 58.7 & \\
\texttt{c\_hem\_s} & ideal & 64.0 & \\
\texttt{c\_k0\_mem} & ideal & 44.7 & 记忆类失败率 \%（多摘+少摘停止）\\
\texttt{c\_k4\_mem} & ideal & 30.7 & \\
\texttt{c\_k8\_mem} & ideal & 21.3 & \\
\texttt{c\_hem\_mem} & ideal & 16.0 & \\
\texttt{c\_k0\_op} & ideal & 16.7 & 操作类失败率 \%（掉果+操作失败）\\
\texttt{c\_k4\_op} & ideal & 17.3 & \\
\texttt{c\_k8\_op} & ideal & 20.0 & \\
\texttt{c\_hem\_op} & ideal & 20.0 & \\
\texttt{c\_k0\_err} & ideal & 0.52 & 平均计数误差 |入篮数−N|（颗）\\
\texttt{c\_k8\_err} & ideal & 0.23 & \\
\texttt{c\_hem\_err} & ideal & 0.17 & \\
\texttt{c\_k4\_err} & ideal & 0.34 & \\
\texttt{c\_d4} & ideal & +13.3 & K4−K0 百分点\\
\texttt{c\_d4\_ci} & ideal & [4.7, 22.0] & \\
\texttt{c\_d8} & ideal & +20.0 & K8−K0\\
\texttt{c\_d8\_ci} & ideal & [11.3, 28.7] & \\
\texttt{c\_dh} & ideal & +25.3 & HEM−K0\\
\texttt{c\_dh\_ci} & ideal & [16.7, 34.0] & \\
\texttt{c\_dh\_p} & ideal & \$<\$0.001 & Holm 校正 p（主假设 H1）\\
\texttt{c\_dh8} & ideal & +5.3 & HEM−K8\\
\texttt{c\_dh8\_ci} & ideal & [\$-\$1.3, 12.0] & \\
\texttt{c\_dh\_noguard} & ideal & +22.7 & 关闭安全包络的敏感性分析 HEM−K0\\
\texttt{c\_dh\_noguard\_ci} & ideal & [12.0, 33.3] & \\
\texttt{c\_clip} & ideal & 22.1 & 闭环控制步裁剪率 \%（各条件接近）\\
\texttt{cN\_k0} & ideal & 56.3, 37.3, 23.5 & K0 成功率 \%，N=2 / N=3 / N=4–5（n=48/51/51）\\
\texttt{cN\_k8} & ideal & 68.8, 58.8, 49.0 & \\
\texttt{cN\_hem} & ideal & 70.8, 62.7, 58.8 & \\
\texttt{cN\_gain\_h0} & ideal & +14.5, +25.4, +35.3 & HEM−K0 按配额（由 cN 推出）\\
\texttt{cN\_gain\_h8} & ideal & +2.0, +3.9, +9.8 & HEM−K8 按配额（由 cN 推出）\\
\texttt{o\_k0} & ideal & 92, 22, 14, 8, 14 & \\
\texttt{o\_k4} & ideal & 94, 20, 13, 9, 14 & \\
\texttt{o\_k8} & ideal & 91, 21, 14, 9, 15 & \\
\texttt{o\_hem} & ideal & 93, 21, 13, 9, 14 & \\
\texttt{o\_k0\_s} & ideal & 61.3 & \\
\texttt{o\_k4\_s} & ideal & 62.7 & \\
\texttt{o\_k8\_s} & ideal & 60.7 & \\
\texttt{o\_hem\_s} & ideal & 62.0 & \\
\texttt{o\_k0\_mem} & ideal & 24.0 & \\
\texttt{o\_k4\_mem} & ideal & 22.0 & \\
\texttt{o\_k8\_mem} & ideal & 23.3 & \\
\texttt{o\_hem\_mem} & ideal & 22.7 & \\
\texttt{o\_d4} & ideal & +1.3 & \\
\texttt{o\_d4\_ci} & ideal & [\$-\$6.0, 8.7] & \\
\texttt{o\_d8} & ideal & \$-\$0.7 & \\
\texttt{o\_d8\_ci} & ideal & [\$-\$8.0, 6.7] & \\
\texttt{o\_dh} & ideal & +0.7 & \\
\texttt{o\_dh\_ci} & ideal & [\$-\$6.7, 8.0] & \\
\texttt{o\_dh\_tost} & ideal & 0.004 & TOST 等效检验 p（±10 pp）\\
\texttt{pc\_cnt\_a} & ideal & 44.7 & 计数 v2 记忆类失败 K0 \%\\
\texttt{pc\_cnt\_b} & ideal & 74.0 & recent16\\
\texttt{pc\_cnt\_d} & ideal & +29.3 & \\
\texttt{pc\_cnt\_ci} & ideal & [19.3, 39.3] & \\
\texttt{pc\_occ\_a} & ideal & 24.0 & \\
\texttt{pc\_occ\_b} & ideal & 61.3 & \\
\texttt{pc\_occ\_d} & ideal & +37.3 & \\
\texttt{pc\_occ\_ci} & ideal & [26.7, 48.0] & \\
\texttt{row\_events} & ideal & 8.6 & 每回合平均采摘次数\\
\texttt{row\_steps} & ideal & 2140 & 每回合平均步数\\
\texttt{aud\_row\_o1} & ideal & 6 & 捷径审计 O1\\
\texttt{aud\_row\_o2} & ideal & 10 & O2\\
\texttt{aud\_row\_o3} & ideal & 94 & O3\\
\texttt{exp\_row} & ideal & 91.0 & 专家逐尝试成功率 \%\\
\texttt{l\_k0\_s} & ideal & 12.7 & K0 成功率 \%\\
\texttt{l\_k8\_s} & ideal & 30.0 & 平坦 PEKM k=8\\
\texttt{l\_hem\_s} & ideal & 52.7 & HEM\\
\texttt{l\_err\_k0} & ideal & 2.41 & 平均计数误差（颗）\\
\texttt{l\_err\_k8} & ideal & 1.37 & \\
\texttt{l\_err\_hem} & ideal & 0.62 & \\
\texttt{l\_k0\_mem} & ideal & 63.3 & 记忆类失败率 \%\\
\texttt{l\_k8\_mem} & ideal & 45.3 & \\
\texttt{l\_hem\_mem} & ideal & 22.0 & \\
\texttt{l\_k0\_op} & ideal & 24.0 & 操作类失败率 \%\\
\texttt{l\_k8\_op} & ideal & 24.7 & \\
\texttt{l\_hem\_op} & ideal & 25.3 & \\
\texttt{l\_dh0} & ideal & +40.0 & HEM−K0\\
\texttt{l\_dh0\_ci} & ideal & [30.7, 49.3] & \\
\texttt{l\_dh8} & ideal & +22.7 & HEM−K8（主假设 H3）\\
\texttt{l\_dh8\_ci} & ideal & [13.3, 32.0] & \\
\texttt{l\_p} & ideal & \$<\$0.001 & H2、H3 的 Holm 校正 p\\
\texttt{lN\_k0} & ideal & 20.0, 10.2, 7.8 & N=6–7 / 8–9 / 10–12（n=50/49/51）\\
\texttt{lN\_k8} & ideal & 44.0, 26.5, 19.6 & \\
\texttt{lN\_hem} & ideal & 60.0, 51.0, 47.1 & \\
\texttt{ab\_oracle} & ideal & 65.3 & 仿真物理事件（上界）\\
\texttt{ab\_pedet} & ideal & 64.0 & PEDet（默认）\\
\texttt{ab\_gap} & ideal & 1.3 & 上界与 PEDet 之差（百分点）\\
\texttt{ab\_geo} & ideal & 58.0 & 几何距离事件（真值距离）\\
\texttt{ab\_fdm} & ideal & 55.3 & SAM-FDM\\
\texttt{ab\_grip} & ideal & 56.7 & 夹爪启发式\\
\texttt{ab\_light\_k0} & ideal & 35.3 & 光照扰动：K0\\
\texttt{ab\_light\_pedet} & ideal & 60.7 & 光照扰动：HEM+PEDet\\
\texttt{ab\_light\_fdm} & ideal & 40.7 & 光照扰动：HEM+SAM-FDM\\
\texttt{ab\_shift} & ideal & 41.3 & 关键帧错后 24 帧\\
\texttt{ab\_random} & ideal & 40.0 & 预留位置随机选取\\
\texttt{ab\_contact} & ideal & 55.3 & 只保留接触帧\\
\texttt{ab\_release} & ideal & 60.0 & 只保留脱离帧\\
\texttt{ab\_untrained} & ideal & 36.0 & K0 策略推理时直接换成 HEM（训练未见）\\
\texttt{ab\_noept} & ideal & 59.3 & 去掉事件扰动训练（需重训）\\
\texttt{ab\_noise\_ept} & ideal & 60.7 & 注入 10\% 漏检：有 EPT\\
\texttt{ab\_noise\_noept} & ideal & 50.0 & 注入 10\% 漏检：无 EPT\\
\texttt{al\_full} & ideal & 52.7 & 完整 HEM\\
\texttt{al\_nosum} & ideal & 26.7 & 去掉情景摘要层（容量 2 次，低于 K8）\\
\texttt{al\_single} & ideal & 38.7 & 摘要层用单帧代替拼图（容量 5 次）\\
\texttt{al\_mosaic9} & ideal & 49.3 & 3×3 拼图（容量 29 次，单格分辨率更低）\\
\texttt{al\_noemb} & ideal & 45.3 & 去掉类型-时间嵌入\\
\texttt{al\_noept} & ideal & 46.0 & 去掉事件扰动训练\\
\texttt{rb\_arm} & todo & RealMan RM75-B 七自由度机械臂 & 按实际设备确认\\
\texttt{rb\_grip} & todo & EG2 两指电动夹爪（或同类） & \\
\texttt{rb\_cam\_f} & todo & Intel RealSense D435i & 前置相机\\
\texttt{rb\_cam\_w} & todo & Intel RealSense D405 & 腕部相机\\
\texttt{rb\_cam\_3} & todo & 60 fps 第三视角相机 & 仅用于评测标注\\
\texttt{rb\_gpu} & todo & NVIDIA RTX 4090 工作站 & \\
\texttt{rb\_hz} & todo & 20 & 控制频率 Hz，应与仿真一致\\
\texttt{rb\_fruit} & todo & 番茄（横径 5--6 cm） & 真实果实品种与尺寸\\
\texttt{rb\_force} & todo & 3--5 & 果梗夹脱离力 N\\
\texttt{rb\_demo} & todo & 120 & 每个任务的遥操作示范回合数\\
\texttt{rb\_hsl\_ep} & ideal & 80 & 用于事后确认自训练的无标注真机回合数\\
\texttt{rb\_sup\_ep} & ideal & 20 & 人工标注参照组的回合数\\
\texttt{rb\_det\_ep} & ideal & 40 & 真机检测评测回合数\\
\texttt{rb\_det\_ev} & ideal & 118 & 真机检测评测事件数\\
\texttt{rb\_annot} & ideal & 1 & 两名标注者脱离帧差异中位数（帧）\\
\texttt{rb\_annot2} & ideal & 96.8 & 两名标注者差异不超过 2 帧的事件比例 \%\\
\texttt{hsl\_keep} & ideal & 83.5 & 通过一致性筛选、被用作伪标签的持物段比例 \%\\
\texttt{hsl\_prec} & ideal & 0.991 & 伪标签相对人工标注的事件精确率\\
\texttt{rd\_grip} & ideal & 0.872, 0.941, 0.905 & \\
\texttt{rd\_fdm} & ideal & 0.861, 0.912, 0.886 & \\
\texttt{rd\_zero} & ideal & 0.842, 0.873, 0.857 & PEDet 仅仿真训练\\
\texttt{rd\_hsl} & ideal & 0.957, 0.946, 0.951 & PEDet + 事后确认自训练（无人工标注）\\
\texttt{rd\_sup} & ideal & 0.963, 0.958, 0.960 & PEDet + 20 回合人工标注（参照）\\
\texttt{rd\_tr} & ideal & 1 & 真机脱离关键帧误差中位数（帧）\\
\texttt{fig\_hsl\_pedet} & ideal & 0.857, 0.902, 0.927, 0.944, 0.951 & \\
\texttt{fig\_sup\_ref} & ideal & 0.96, 0.96, 0.96, 0.96, 0.96 & \\
\texttt{fig\_hsl\_fdm} & ideal & 0.886, 0.886, 0.886, 0.886, 0.886 & \\
\texttt{fig\_hsl\_grip} & ideal & 0.905, 0.905, 0.905, 0.905, 0.905 & \\
\texttt{rc\_k0} & ideal & 10, 11, 5, 2, 2 & 真机计数：K0\\
\texttt{rc\_k8} & ideal & 17, 5, 3, 2, 3 & 真机计数：平坦 PEKM k=8\\
\texttt{rc\_hem} & ideal & 20, 2, 3, 2, 3 & 真机计数：HEM\\
\texttt{rc\_k0\_s} & ideal & 33.3 & \\
\texttt{rc\_k8\_s} & ideal & 56.7 & \\
\texttt{rc\_hem\_s} & ideal & 66.7 & \\
\texttt{rc\_d} & ideal & +33.3 & HEM−K0\\
\texttt{rc\_p} & ideal & 0.004 & 布局配对符号翻转检验 p\\
\texttt{rc\_d8} & ideal & +10.0 & HEM−K8\\
\texttt{rc\_p8} & ideal & 0.38 & \\
\texttt{rc\_err\_k0} & ideal & 0.63 & 平均计数误差\\
\texttt{rc\_err\_hem} & ideal & 0.20 & \\
\texttt{ro\_k0\_s} & ideal & 53.3 & 真机遮挡：K0 16/30\\
\texttt{ro\_hem\_s} & ideal & 56.7 & 真机遮挡：HEM 17/30\\
\texttt{rl\_pedet\_s} & ideal & 60.0 & 真机暖光扰动计数：HEM+PEDet 18/30\\
\texttt{rl\_fdm\_s} & ideal & 36.7 & 真机暖光扰动计数：HEM+SAM-FDM 11/30\\
\texttt{gh\_k0} & ideal & 5/15 & 温室盆栽番茄：K0\\
\texttt{gh\_hem} & ideal & 10/15 & 温室盆栽番茄：HEM\\
\texttt{rt\_k0} & ideal & 112 & 真机单回合平均用时 s\\
\texttt{rt\_hem} & ideal & 115 & \\
\texttt{rt\_vla} & ideal & 96 & 真机 VLA 单次推理 ms\\
\texttt{rt\_pedet} & ideal & 2.4 & 真机 PEDet 单帧 ms\\
\texttt{sim\_vla\_ms} & todo & 【待测】 & 仿真服务端 VLA 单次推理 ms（L20）\\
\bottomrule
\end{longtable}

}
